\section{Introduction}
\label{sec:introduction}

Bloch's law predicts that the spontaneous magnetization of a
three-dimensional ferromagnet falls below its fully polarized value by
an amount proportional to $T^{3/2}$ at low temperature. We prove this
law for the quantum Heisenberg ferromagnet at every fixed spin, allowing
any nonnegative symmetric finite-range interaction whose support
generates the lattice. The interaction may be spatially anisotropic and need not
contain the ordinary nearest-neighbor bonds. Its quadratic dispersion
determines the leading coefficient.

\subsection{The model and the result}

Fix $S\in\{\tfrac12,1,\tfrac32,\ldots\}$. The spin space at each
site is $\mathfrak h_S=\C^{2S+1}$, with orthonormal basis
$(e_m)_{m=-S}^S$ and operators
\[
 S^ze_m=me_m,\qquad
 S^+e_m=\sqrt{S(S+1)-m(m+1)}\,e_{m+1}\quad(m<S),\qquad S^+e_S=0.
\]
Set $S^-=(S^+)^*$, $S^x=(S^++S^-)/2$, and
$S^y=(S^+-S^-)/(2i)$. A subscript specifies the site where an operator
acts. Let $J:\Z^3\to[0,\infty)$ have finite support and satisfy
\begin{equation}
 J(0)=0,\qquad J(z)=J(-z),\qquad
 \langle z:J(z)>0\rangle_{\Z}=\Z^3.
 \label{eq:interaction}
\end{equation}
On $\Lambda_N=[-N,N]^3\cap\Z^3$, with free boundary, define
\begin{equation}
 H^J_{S,N}=-\frac12\sum_{x,y\in\Lambda_N}
               J(y-x)\boldsymbol S_x\cdot\boldsymbol S_y,
 \qquad M_{S,N}=\sum_{x\in\Lambda_N}S_x^z.
 \label{eq:hamiltonian}
\end{equation}
The sum counts each unordered bond twice; the factor $1/2$ removes
that duplication. We set the lattice spacing, Boltzmann constant, and
$\hbar$ equal to one, so $T=\beta^{-1}$.

For $\beta>0$ and $h\in\R$, put
\begin{equation}
 p_\beta(h)=\lim_{N\to\infty}\frac1{\beta|\Lambda_N|}
       \log\Tr e^{-\beta(H^J_{S,N}-hM_{S,N})},
 \qquad m_{S,J}(\beta)=\partial_h^+p_\beta(0).
 \label{eq:pressure-magnetization}
\end{equation}
The finite-range pressure limit exists and is convex and $S$-Lipschitz;
Section~\ref{sec:pressure} recalls the proof. Rotation by $\pi$ about
the $x$-axis makes the pressure even, so $m_{S,J}(\beta)\in[0,S]$.
The derivative selects the magnetization in a field tending to zero
through positive values after the thermodynamic limit.

Define the positive definite matrix
\begin{equation}
 D_{S,J}=\frac S2\sum_{z\in\Z^3}J(z)zz^{\mathsf T}.
 \label{eq:diffusion-matrix}
\end{equation}
It is the quadratic coefficient of the one-magnon dispersion:
\begin{equation}
 \varepsilon_{S,J}(k)=S\sum_zJ(z)(1-\cos(k\cdot z))
       =k^{\mathsf T}D_{S,J}k+O(|k|^4).
 \label{eq:dispersion}
\end{equation}
Positive definiteness follows because the interaction support spans
$\R^3$.

\begin{theorem}[Bloch's law]
\label{thm:main}
For every fixed $S\in\{\tfrac12,1,\tfrac32,\ldots\}$ and every
interaction satisfying \eqref{eq:interaction},
\begin{equation}
 \lim_{\beta\to\infty}\beta^{3/2}\bigl(S-m_{S,J}(\beta)\bigr)
   =\frac{\zeta(3/2)}{8\pi^{3/2}\sqrt{\det D_{S,J}}}.
 \label{eq:bloch-law}
\end{equation}
\end{theorem}

For $J(\pm e_j)=1$ and all other couplings zero, $D_{S,J}=SI$;
the theorem gives the usual nearest-neighbor coefficient
$\zeta(3/2)/(8\pi^{3/2}S^{3/2})$. The order of limits is part of
the assertion: first volume tends to infinity at fixed $\beta,h$,
then $h\downarrow0$, and finally $\beta\to\infty$, with $S,J$ fixed.

\subsection{Historical context and dependencies}

Bloch derived the $T^{3/2}$ reduction from long-wavelength spin waves
\cite{Bloch1930}. The bosonic representation of Holstein and Primakoff
\cite{HolsteinPrimakoff1940} and Dyson's analysis of spin-wave interactions
\cite{Dyson1956a,Dyson1956b} developed the low-temperature expansion.
In particular, Dyson identified the interaction contribution to the
magnetization at order $T^4$; his estimates of the remaining higher-order
terms were not a rigorous fixed-spin proof of the asymptotic law
\cite[Section~9]{Dyson1956b}. The leading coefficient in
\eqref{eq:bloch-law} is obtained from the low-temperature density of ideal magnons with quadratic
energy $k^{\mathsf T}D_{S,J}k$. The issue is to justify this coefficient
for the interacting spin system and the field derivative in
\eqref{eq:pressure-magnetization}.

Rigorous ordering theory distinguishes sharply between dimensions and
between classical and quantum models. Mermin and Wagner excluded
positive-temperature spontaneous magnetic order for the one- and
two-dimensional isotropic Heisenberg models under their finite-range
hypotheses \cite{MerminWagner1966}. Infrared bounds proved symmetry
breaking for classical models \cite{FSS1976}, and Dyson, Lieb, and Simon
established ordering for classes of quantum antiferromagnets
\cite{DLS1978}. These developments left the standard three-dimensional
quantum ferromagnet as a separate problem, recorded by Lieb
\cite{Lieb1999} and still discussed as open in recent accounts
\cite{Seiringer2025,BjornbergRyan2026}. The distinct companion manuscript
\cite[Theorem~1.1]{Ordering} proves low-temperature magnetized equilibrium states
for the nearest-neighbor quantum ferromagnet. We use its finite-set
algebra and deterministic lattice geometry, with their precise
hypotheses stated at the points of use.

The free energy has a different rigorous history. Conlon and Solovej
\cite{ConlonSolovej1991} and T\'oth \cite{Toth1993} obtained bounds for
spin $1/2$. Correggi, Giuliani, and Seiringer proved the leading
low-temperature free-energy asymptotic for the nearest-neighbor model
on $\Z^3$ at every fixed spin \cite[Theorem~2.1]{CGS2015}. A zero-field free-energy estimate
does not by itself control the one-sided field derivative that defines
spontaneous magnetization. Our proof therefore estimates the observable
through a representation that remains uniform under arbitrary extra
up-spin constraints, and only then passes to pressure increments.

The probabilistic representation belongs to the random-interchange
approach of T\'oth \cite{Toth1993}; the symmetric-slot realization of
higher spin is described by Nachtergaele \cite{Nachtergaele1994}.
The pin argument uses the stable-polynomial approach to negative
dependence developed by Borcea, Br\"and\'en, and Liggett \cite{BBL2009}. Diffusive smoothing uses the
energy method of Nash \cite{Nash1958}. The companion \cite{Ordering}
provides the particular stable exchange symbols, presence inequality,
local lattice inequality, and routed unit flows used below. We prove
here the finite-range pin bootstrap, the quantitative long-return
bound, and the comparison that recovers the full dispersion.

\subsection{How returns determine the coefficient}

Represent each spin by $\ell=2S$ spin-$1/2$ \emph{slots}. Independent
Poisson exchanges between slots, followed by a uniform permutation of
the slots at each site, form directed cycles over a time period of
length $\beta$. A cycle has a single up or down color. Constraining
all cycles through the boundary and through specified slot-time points
to be up defines a \emph{pin law}. If $E$ is its pin set and $i$ is
an unpinned slot-time point, a change of weights gives
\[
 d_E(i)=\frac{r(E,\{i\})}{1+r(E,\{i\})},
\]
where $d_E(i)$ is its down probability. Under the law with $i$ also
pinned, $r(E,\{i\})$ is the probability that the directed line from
$i$ returns to $i$ at positive elapsed time before meeting $E$.
Section~\ref{sec:pins} makes this identity precise and represents
$r$ by a walk on the slots left after exposing the cycles missing the
augmented pin set. The exposed environment repeats over periods, while
the walk uses fresh transition randomness on every traversal.

There are two geometric tasks. For rough estimates, choose three
linearly independent interaction vectors. They generate a sublattice
of finite index; on each of its cosets, the chosen bonds form a copy
of the ordinary cubic lattice. Section~\ref{sec:geometry} transfers
local Nash and flow estimates to these copies and compares their
energies with the full interaction energy. Section~\ref{sec:bootstrap}
then proves a sparse-pin estimate by downward induction on the pin
set. This induction supplies enough regular time for diffusion without
assuming a finite-range ordering theorem.

Sharp asymptotics require all interaction bonds. The ideal walk has
site jump rates $SJ(z)$, and its heat kernel $p_t^{S,J}$ satisfies
\[
 p_t^{S,J}(x,x)\sim\kappa t^{-3/2},\qquad
 \kappa=(4\pi)^{-3/2}(\det D_{S,J})^{-1/2}.
\]
After a full traversal the slot index is uniform, so the ideal
$n$-period return probability to a specified slot is
$\ell^{-1}p_{n\beta}^{S,J}(x,x)$. Summing over the $\ell$ slots at a
site cancels this factor, leaving the finite sums
$\kappa\sum_{n\le K}n^{-3/2}$ that approach the Bloch coefficient.
The selected sublattice is used only for estimates; the full jump
kernel determines this coefficient.

The main analytic difficulty is controlling arbitrarily late returns
in the same repeated environment. Section~\ref{sec:returns} truncates
a quadratic hitting-probability identity after $K$ periods. A mass
already spread for time comparable to $K\beta$ has a flow to infinity
with energy at most $C(K\beta)^{-1/2}$. Averaging over one further
period gives a remainder bounded by
$CK^{-1/2}\beta^{-3/2}+o_K(\beta^{-3/2})$, with $C$ independent of
$K$. Its $K=1$ case first improves the exposed-slot density to
$O(\beta^{-3/2})$.

This improved density permits the comparison in
Section~\ref{sec:ideal}. The ideal path rarely encounters an exposed
strand over a fixed number of periods. Splitting a return at its
midpoint combines that small encounter probability with diffusive
smoothing, giving an error $o_K(\beta^{-3/2})$ in each return entry.
The comparison retains the same exposed environment in both halves;
independence of their environments is unnecessary. The finite-return
comparison and the long-return remainder give matching upper and
lower down-probability bounds, with the lower bound requiring a
pin-free neighborhood of the starting point.

Finally, Section~\ref{sec:pressure} expresses a positive field as a
mixture of extra pins. The probability of selecting any given slot
is at most $1-e^{-\beta h}$, so the required neighborhood is free of
extra pins with probability tending to one as $h\downarrow0$ at fixed
$\beta,K$. Integrating the finite-volume derivative bounds before
taking the thermodynamic limit proves \eqref{eq:bloch-law} in the
stated order. The uniform pin estimates, the finite-index geometric
reduction, and the energy bound for a spread source are the reusable
parts of the argument.
